% Generated draft: requires mathematical and source review.
\documentclass[11pt]{article}
\usepackage[margin=25mm]{geometry}
\usepackage{fontspec}
\setmainfont{Latin Modern Roman}
\usepackage{amsmath,amssymb,mathtools,unicode-math}
\setmathfont{Latin Modern Math}
\usepackage{enumitem,needspace,xurl,hyperref}
\setlength{\parindent}{0pt}
\setlength{\parskip}{5pt}
\setlength{\emergencystretch}{3em}
\newcommand{\N}{\mathbb{N}}
\newcommand{\Z}{\mathbb{Z}}
\newcommand{\Q}{\mathbb{Q}}
\newcommand{\R}{\mathbb{R}}
\newcommand{\C}{\mathbb{C}}
\newcommand{\F}{\mathbb{F}}
\newcommand{\A}{\mathbb{A}}
\newcommand{\PP}{\mathbb{P}}
\newcommand{\E}{\mathbb{E}}
\newcommand{\eps}{\varepsilon}
\newcommand{\dd}{\,\mathrm{d}}
\begin{document}
\section*{UnsolvedMath ID: 20000373}
\textbf{Source problem code:} catalogue entry
\subsection{Definitions and mathematical statement}
Let $G\hookrightarrow\mathrm{SL}_{n,k}$ be a finite group scheme over a number field $k$, and put $X=\mathrm{SL}_{n,k}/G$. In characteristic zero such finite group schemes are étale, but need not be constant: $G(\overline k)$ carries a possibly nontrivial Galois action. Write $H^1(k,G)$ for continuous nonabelian Galois cohomology, a pointed set, and $\operatorname{loc}_{\R}$ for localization at the real places. All density statements below use the ordinary real topology and the diagonal image of rational points.

\textbf{Original Problem 5, Harari reporting a question of Borovoi.} For $k=\Q$ and arbitrary such $G$, is $X(\Q)$ dense in $X(\R)$? Equivalently, is
\[H^1(\operatorname{Gal}(\overline\Q/\Q),G(\overline\Q))\longrightarrow H^1(\operatorname{Gal}(\C/\R),G(\C))\]
surjective? The source says this is known for constant $G$. It also asks: if $G$ is constant over a number field $K$ with $r\ge2$ real places $v_1,\ldots,v_r$, is
\[X(K)\text{ dense in }\prod_{j=1}^rX(K_{v_j})?\]
These are real-place instances of the broader weak-approximation question; finite places are not silently included in the stated products. Lucchini Arteche comments that the algebraic Brauer--Manin obstruction says nothing for this problem. Skorobogatov attributes the abelian case to Borovoi.

\textbf{Attached unreviewed 2026 status claims.} The aid says a 2026 preprint resolves the constant-group several-real-places variant affirmatively, and another 2026 preprint together with the known real-place Brauer result covers solvable finite group schemes. It describes the arbitrary nonconstant nonsolvable case as open. Neither preprint nor the cited Brauer result is bibliographically identified, so these are dated supplied claims, not a new literature assessment.

The aid also proposes an odd-kernel reduction. Let
\[1\longrightarrow N\longrightarrow A\xrightarrow{p}Q\longrightarrow1\]
be an exact sequence of finite $k$-group schemes with $N$ of odd geometric order. It asserts that for every real place $v$,
\[p_{*,v}:H^1(k_v,A)\xrightarrow{\sim}H^1(k_v,Q)\]
is a bijection of pointed sets. Define real approximation for a group scheme $H$ as surjectivity of $H^1(k,H)\to\prod_{v\text{ real}}H^1(k_v,H)$. The claimed exact global criterion is
\[A\text{ has real approximation}\quad\Longleftrightarrow\quad\operatorname{loc}_{\R}\bigl(p_*H^1(k,A)\bigr)=\prod_{v\text{ real}}H^1(k_v,Q).\]
It is the image of globally liftable $Q$-classes that occurs here, not automatically all of $H^1(k,Q)$. If the extension has a $k$-group section, equivalently a Galois-equivariant group section on geometric points, the aid concludes that real approximation for $A$ is equivalent to that for $Q$.

Finally, the aid identifies the connected components of $(\mathrm{SL}_n/A)(\R)$ with $H^1(\R,A)$ and reports the equivalence between localization surjectivity and density on products of real loci. These identifications use the finite embedding, the usual $\mathrm{SL}_n$ cohomology and approximation framework, and require review along with the reduction.
\subsection{Short English statement}
The source asks whether real local classes for finite stabilizers come from global points, first over $\Q$ and then simultaneously at several real places. The aid reports dated progress and an odd-order-kernel reduction, carefully retaining the global liftability condition and the extra splitting hypothesis needed to reduce directly to the quotient group.
\subsection{Sources}
\begin{itemize}
\item Local UnsolvedMath registry, permanent ID 20000373, source code catalogue entry; \texttt{aim-algebraic-number-theory-notes.json}, notes[24], tag \texttt{problem}.
\item AIM \emph{Rational and integral points on higher-dimensional varieties}, open problem session, May 28, 2014, Problem 5: \url{https://aimath.org/pastworkshops/ratlhigherdimvarproblems.pdf}. Harari attributes the question to Borovoi; comments are attributed to Lucchini Arteche and Skorobogatov. The supplied verification describes restored overlines and product indices; no PDF was consulted here.
\item Attached machine-generated research attempt 2 mentions two 2026 preprints and a real-place Brauer result without authors, titles or identifiers. Result \texttt{partial\_result}, status \texttt{partially\_solved}, rigor label \texttt{proved}, novelty confidence low. It is not peer reviewed and requires independent verification before citation; original literature field empty.
\end{itemize}
\end{document}
